\subsection{Object Persistence}\label{sec:status-records-coherence:object-persistence}

\begin{theorem}[Object Persistence]\label{thm:F19}
Let \(H_i,H_j\) be carriers with access quotients
\(q_i:H_i\to Q_i\) and \(q_j:H_j\to Q_j\), and let
\(\gamma:H_i\to H_j\) be a declared transport.  Assume \(q_i\) and
\(q_j\) are surjective.  Let
\(\operatorname{viable}:Q_j\to\mathrm{Prop}\) be a viability
predicate, \(C:H_i\times H_j\to\mathrm{Prop}\) a relation predicate,
and let
\[
  \bar\gamma:Q_i\to Q_j,\qquad
  \bar\eta:Q_j\to Q_k,\qquad
  \operatorname{compBar}:Q_i\to Q_k
\]
be arbitrary holonomy data (no compatibility between
\(\operatorname{compBar}\) and \(\bar\eta\circ\bar\gamma\) is
assumed).  Then
\begin{align*}
  &\operatorname{Persists}(q_i,\gamma,q_j)
    \iff
    \operatorname{PersistenceObstructionEmpty}(q_i,\gamma,q_j),\\
  &\operatorname{SplitObstructionEmpty}(q_i,\gamma,q_j)
    \iff
    \operatorname{Persists}(q_i,\gamma,q_j),\\
  &\operatorname{ViablePersistence}(\gamma,q_j,\operatorname{viable})
    \iff
    \operatorname{DissolutionObstructionEmpty}
      (\gamma,q_j,\operatorname{viable}),\\
  &\operatorname{RelationValuedPersistence}(q_i,q_j,C)
    \iff
    \operatorname{RelationPersistenceObstructionEmpty}(q_i,q_j,C),\\
  &\operatorname{IdentityHolonomyCoherent}
      (\bar\gamma,\bar\eta,\operatorname{compBar})\\
  &\qquad{}\iff
    \operatorname{IdentityHolonomyDefectEmpty}
      (\bar\gamma,\bar\eta,\operatorname{compBar}),
\end{align*}
together with the comparison-status partition
\[
  \operatorname{ExactlyOnePersistenceComparisonStatus}
    (q_i,\gamma,q_j),
\]
and the retention and carrying conditions
\begin{align*}
  &\operatorname{RetainsOldIdentity}
      \bigl(q_i,\operatorname{persistenceExtension}(q_i,\gamma,q_j)\bigr),\\
  &\operatorname{CarriesLaterIdentity}
      \bigl(\operatorname{persistenceExtension}(q_i,\gamma,q_j),
        \gamma,q_j\bigr).
\end{align*}
\end{theorem}
\footnote{\raggedright Verified in Lean as
\texttt{StatusRecordsCoherence.ObjectPersistence.\allowbreak{}object\_persistence};
see \Cref{app:formalization}.}
\begin{proof}
Assume the surjectivity of \(q_i\) and \(q_j\), the predicates
\(\operatorname{viable}\) and \(C\), and the holonomy data
\(\bar\gamma,\bar\eta,\operatorname{compBar}\).  We prove the eight
conjuncts in turn.

\emph{(1) Persists iff the persistence obstruction is empty.}
The persistence obstruction is
\[
  \Delta_{\mathrm{pers}}(q_i,\gamma,q_j)
  =
  \{(h,h')\in H_i\times H_i:
      q_i(h)=q_i(h')\ \text{and}\
      q_j(\gamma h)\ne q_j(\gamma h')\}.
\]
Its emptiness is exactly the quotient-descent condition for
\(q_j\circ\gamma:H_i\to Q_j\) through \(q_i\).  Since \(q_i\) is
surjective, this is equivalent to the existence of a descended map
\(\bar\gamma:Q_i\to Q_j\) satisfying
\[
  \bar\gamma\circ q_i=q_j\circ\gamma,
\]
which is \(\operatorname{Persists}(q_i,\gamma,q_j)\).

\emph{(2) Split-obstruction empty iff persists.}
The split obstruction is the same pair set written in the comparison
vocabulary:
\[
  \Delta_{\mathrm{split}}(q_i,\gamma,q_j)
  =
  \Delta_{\mathrm{pers}}(q_i,\gamma,q_j).
\]
Thus \(\Delta_{\mathrm{split}}=\emptyset\) iff the persistence
obstruction is empty, and by (1) this is equivalent to persistence.

\emph{(3) Exactly-one persistence comparison status.}
The comparison statuses are decided by two split-pair defects on the
declared continuation \(\gamma\): the persistence/split obstruction
\(\Delta_{\mathrm{split}}\) of case~(1), and the merge obstruction
\[
  \Delta_{\mathrm{merge}}(q_i,\gamma,q_j)
  =
  \{(h,h')\in H_i\times H_i:
      q_j(\gamma h)=q_j(\gamma h')\ \text{and}\ q_i(h)\ne q_i(h')\}.
\]
The four statuses partition the configurations by the joint
\(\Delta_{\mathrm{split}}\)\textendash{}\(\Delta_{\mathrm{merge}}\)
emptiness pattern:
\[
\begin{array}{c|c|c}
\text{status} & \Delta_{\mathrm{split}}=\emptyset
  & \Delta_{\mathrm{merge}}=\emptyset\\
\hline
\text{exact endurance} & \text{yes} & \text{yes}\\
\text{persistent with merge} & \text{yes} & \text{no}\\
\text{identity branching} & \text{no} & \text{yes}\\
\text{identity rewrite} & \text{no} & \text{no}
\end{array}
\]
Every configuration of the two Boolean defect tests falls into
exactly one row.  Existence is by case enumeration; exclusivity is by
definitional disjointness of the four status predicates.  The
dissolution and relation-valued obstructions of cases~(4) and~(5) are
separate typed gates and do not enter this partition.

\emph{(4) Viable persistence iff the dissolution obstruction is
empty.}
The dissolution obstruction is
\[
  \Delta_{\mathrm{diss}}(\gamma,q_j,\operatorname{viable})
  =
  \{h\in H_i:
      \neg\,\operatorname{viable}(q_j(\gamma h))\}.
\]
Viable persistence says that every transported \(q_j\)-class lies in
the viable region:
\[
  \forall h\in H_i,\quad
  \operatorname{viable}(q_j(\gamma h)).
\]
This is precisely \(\Delta_{\mathrm{diss}}=\emptyset\).

\emph{(5) Relation-valued persistence iff its obstruction is empty.}
The relation-persistence obstruction is
\[
  \Delta_C(q_i,q_j,C)
  =
  \{(h,h')\in H_i\times H_i:
    q_i(h)=q_i(h')\ \text{and}\
    C(h,\gamma h)\not\iff C(h,\gamma h')\}.
\]
Emptiness of \(\Delta_C\) says that the truth of the relation is
independent of the chosen representative of a \(q_i\)-class.  Hence
\(C\) descends to a well-defined relation on
\(Q_i\times Q_j\), which is exactly relation-valued persistence.

\emph{(6) Identity holonomy coherent iff the holonomy defect is
empty.}
The holonomy defect is
\[
  \Delta_{\mathrm{hol}}(\bar\gamma,\bar\eta,\operatorname{compBar})
  =
  \{a\in Q_i:
      \operatorname{compBar}(a)\ne\bar\eta(\bar\gamma(a))\}.
\]
Thus \(\Delta_{\mathrm{hol}}=\emptyset\) iff
\[
  \forall a\in Q_i,\quad
  \operatorname{compBar}(a)=\bar\eta(\bar\gamma(a)),
\]
which is the pointwise form of
\(\operatorname{compBar}=\bar\eta\circ\bar\gamma\), i.e. identity
holonomy coherence.

\emph{(7) The persistence extension retains the old identity.}
Define the persistence extension as the joint quotient
\[
  \operatorname{persistenceExtension}(q_i,\gamma,q_j)
  :=
  \langle q_i,q_j\circ\gamma\rangle:H_i\to Q_i\times Q_j.
\]
Projection to the first coordinate recovers the old identity:
\[
  \operatorname{pr}_{Q_i}\circ
  \langle q_i,q_j\circ\gamma\rangle=q_i.
\]
Therefore the persistence extension retains \(q_i\).

\emph{(8) The persistence extension carries the later identity.}
Projection to the second coordinate recovers the transported later
identity:
\[
  \operatorname{pr}_{Q_j}\circ
  \langle q_i,q_j\circ\gamma\rangle=q_j\circ\gamma.
\]
Thus the same extension carries \(q_j\) along the declared transport
\(\gamma\).
\end{proof}

The theorem does not classify persistent objects or say which
transports should be declared.

\subsection*{Layer instantiations}\label{layer-inst:F19}

\begin{description}
\item[\emph{Mendelian inheritance (genetics).}]
Through meiotic segregation and fertilisation, parental allele
identity at a locus survives into the offspring's diploid
genotype.  The F19 reading: allele identity persists when no
mutation or non-disjunction breaks the descent; mutation and
non-disjunction events are the persistence obstruction
\citep{Mendel1866VersuchePflanzenHybriden}.

